\section{The balanced semiprime term}\label{sec:k2}

We bound the complete occurrence sum \(B_b(N)\).
The two prime orientations keep their multiplicities.
In this section, fix
\begin{equation}\label{eq:k2-parameters}
\begin{gathered}
 \frac13<b<\frac12,\qquad \varepsilon=\eta_L=10^{-4},\qquad
 u=\frac1{500},\qquad \mu_* =\frac{25}{128},\qquad L=\log N,\\
 0<\sigma\le10^{-10},\qquad \tau=\sigma^9,\qquad \rho=1+\tau,\\
 N'=N(1-L^{-333}),\\
 \mathcal N_K=(1-\varepsilon)U_C=\frac{W(N)}4,\\
 I(b)=\int_b^{1/2}\frac{dt}{t(1-t)}\\
 =\log\frac{1-b}{b}.
\end{gathered}
\end{equation}
The final argument uses \eqref{eq:intro-exponents}.
Every limit below fixes \(b\) and \(\sigma\) first.
For row three, write
\(\ell_0=53/412-\eta_L/2=264897/2060000>0\). This is the fixed floor of
Paper~I, \S3.1.

\subsection{The convolution source}

Define
\begin{equation}\label{eq:k2-source}
\begin{aligned}
 A_b(e)&=\#\{(m,p')\in\mathbb P^2:
 N^b\le m\le\sqrt N,\quad \varepsilon N<mp'=e\le N'\},\\
 X_b&=\sum_e A_b(e),\qquad
 B_b' =\sum_e A_b(e)1_{N-e\in\mathbb P},\\
 E_{\mathrm{top}}&=B_b-B_b',\qquad
 0\le A_b(e)\le2,\qquad
 0\le E_{\mathrm{top}}\le2(NL^{-333}+1).
\end{aligned}
\end{equation}
The last bound counts integers in the removed interval.
A prime square contributes one occurrence.
No new ordering or coprimality restriction is imposed.

For a finite sequence \(F\), put
\begin{equation}\label{eq:k2-natural-kernel}
 \Delta_F(q,a)=\sum_eF(e)
 \left(1_{e\equiv a\ (q)}-\frac{1_{(e,q)=1}}{\varphi(q)}\right),
 \qquad (a,q)=1.
\end{equation}
Both terms use the full support.
For this source, PNT and partial summation give
\begin{equation}\label{eq:k2-source-mass}
 X_b=\frac{(1-\varepsilon)N}{L}I(b)+o_b(N/L).
\end{equation}
This is the calculation of Paper~II,
\cite[(5.44)--(5.51)]{PaperII}.
The moving top removes at most \(2(NL^{-333}+1)\) mass.

The required convolution estimate is
\begin{equation}\label{eq:k2-bv-source}
 \sum_{q\le Q}\max_{(a,q)=1}|\Delta_{A_b}(q,a)|
 \ll_{A,b}NL^{-A},\qquad
 Q\le N^{1/2}L^{-B_{\mathrm{BV}}(A,b)},
\end{equation}
for every fixed \(A>0\).
Its proof keeps the complete occurrence source and the moving top.
Fix a top \(T\in\{N,N'\}\). Before the discrepancy estimate, compare the
exact cell sources \eqref{eq:k2-cell-source} with their upper and lower
convolutions. Each signed child keeps the natural
coprimality comparator \(1_{(e,q)=1}/\varphi(q)\). It is not replaced
after summing the blocks. The step from
\(T=N\) to \(T=N'\) costs the endpoint bands
\eqref{eq:k2-boundary-cost}. It also costs the removed top mass in
\eqref{eq:k2-source}. The actual moving-top application is proved below.

We use the following convolution input.
Let \(a_m,b_n\) have supports \(m\sim M,n\sim Y\).
Put \(x=MY\ge3\), and fix \(\eta,B_{\mathrm{div}}>0\).
Assume
\begin{equation}\label{eq:k2-convolution-sw}
\begin{gathered}
 x^\eta\le M,Y\le x^{1-\eta},\\
 \left|\sum_{\substack{n\equiv l\ (k)\\(n,h)=1}}b_n
 -\frac1{\varphi(k)}\sum_{(n,hk)=1}b_n\right|
 \ll_{A_2}\|b\|_2Y^{1/2}d(h)^{B_{\mathrm{div}}}
                       (\log(2Y))^{-A_2}
\end{gathered}
\end{equation}
uniformly for every fixed \(A_2>0\), \(h,k\ge1\),
and \((l,k)=1\).
Here \(d(n)\) is the divisor function.
Then Fact~\ref{fact:bv-full-support} gives
\begin{equation}\label{eq:k2-convolution}
 \sum_{q\le Q}\max_{(a,q)=1}|\Delta_{a*b}(q,a)|
 \ll_A\|a\|_2\|b\|_2x^{1/2}(\log x)^{-A},
 \qquad Q\le x^{1/2}(\log x)^{-B_{\mathrm{BV}}}.
\end{equation}
The constants depend on all family constants in
\eqref{eq:k2-convolution-sw}.
This specializes Fact~\ref{fact:bv-full-support}. It is
Paper~II, Lemma~5.1 \cite[(5.23)--(5.28)]{PaperII}.
It applies to complete rectangular supports.

We make the interval reduction explicit.
Fix a target \(A>2\), and put \(H_g=A+4\).
Choose
\begin{equation}\label{eq:k2-interval-grid}
\begin{gathered}
 R_g=\left\lceil
 \frac{(1/2-b)L}{\log(1+L^{-H_g})}\right\rceil,\qquad
 \xi=N^{(1/2-b)/R_g},\qquad y_j=N^b\xi^j,\\
 J_j=[y_{j-1},y_j),\qquad Y_j=y_j.
\end{gathered}
\end{equation}
Close the final upper endpoint.
For the fixed top \(T\in\{N,N'\}\), put
\begin{equation}\label{eq:k2-convolution-envelopes}
\begin{aligned}
 b_j(m)&=1_{m\in\mathbb P\cap J_j},\\
 a_{j,T}^+(p')&=1_{p'\in\mathbb P}
       1_{\varepsilon N/Y_j\le p'\le T\xi/Y_j},\\
 a_{j,T}^-(p')&=1_{p'\in\mathbb P}
       1_{\varepsilon N\xi/Y_j<p'<T/Y_j}.
\end{aligned}
\end{equation}
For \(F,G\) on the positive integers write
\((F*G)(e)=\sum_{mn=e}F(m)G(n)\), and put
\(A_{j,T}^{\pm}=a_{j,T}^{\pm}*b_j\). The exact cell source is
\begin{equation}\label{eq:k2-cell-source}
 A_{j,T}(e)=1_{\varepsilon N<e\le T}
       \sum_{\substack{m\in\mathbb P\cap J_j,\ p'\in\mathbb P\\mp'=e}}1.
\end{equation}
The definitions give the pointwise envelopes
\(A_{j,T}^-\le A_{j,T}\le A_{j,T}^+\).
Their difference is supported in multiplicative boundary bands at
\(\varepsilon N\) and \(T\). Put
\(\mathcal E_Q(F)=\sum_{q\le Q}\max_{(a,q)=1}|\Delta_F(q,a)|\).
Let
\(B_T(e)=\sum_j(A_{j,T}^+(e)-A_{j,T}^-(e))\). Then \(B_T\ge0\), with
support in
\([\varepsilon N/\xi,\varepsilon N\xi]\cup[T/\xi,T\xi]\). Each integer
has at most two ordered factorizations into primes. Since
\(1<\xi\le1+L^{-H_g}\) and \(T\in\{N,N'\}\) satisfies \(T\le N\),
\[
 \sum_e B_T(e)\le
 2\{\varepsilon N(\xi-\xi^{-1})+T(\xi-\xi^{-1})+4\}
 \ll NL^{-H_g}+1.
\]
Since \(\xi-\xi^{-1}\le2(\xi-1)\le2L^{-H_g}\), counting integers in the
two complete boundary bands gives the displayed mass bound.
Paper~II, (5.47), gives the aggregate boundary discrepancy
\begin{equation}\label{eq:k2-boundary-cost}
 \mathcal E_Q\!\left(\sum_j(A_{j,T}^+-A_{j,T}^-)\right)
       \ll NL^{-H_g+1}+Q.
\end{equation}
The \(Q\) term includes the bounded-multiplicity changes at exact
endpoints. Split the second factor into dyadic blocks \((M,2M]\):
\[
 a_{j,T,M}^{\pm}(p')=a_{j,T}^{\pm}(p')1_{M<p'\le2M}.
\]
For each block meeting the support,
\(\varepsilon N/(2Y_j)<M\le T\xi/Y_j\), so
\(M\asymp_{\varepsilon}N/Y_j\), with only
\(O_{\varepsilon}(1)\) blocks. Here \(n\sim Z\) means
\(Z<n\le2Z\), as in Section~\ref{sec:classical-inputs}. For the short
factor, the theorem-scale parent is $Y=Y_j/2$; since $\xi<2$ eventually,
$J_j\subset(Y_j/2,Y_j]$ lies in that parent. The norms are
\begin{equation}\label{eq:k2-interval-norms}
 \|a_{j,T,M}^+\|_2\ll\sqrt M,\qquad
 \|a_{j,T,M}^-\|_2\le\|a_{j,T,M}^+\|_2,\qquad
 \|b_j\|_2\ll\sqrt{Y_j}L^{-(H_g+1)/2}.
\end{equation}
The prime number theorem on these intervals also gives a matching
lower bound. It is \(\|b_j\|_2\gg\sqrt{Y_j}L^{-(H_g+1)/2}\) for all sufficiently
large \(N\). This lower bound normalizes the conversion of the absolute
prime Siegel--Walfisz error to
\eqref{eq:k2-convolution-sw}.
The prime-interval form of \eqref{eq:k2-convolution-sw} holds uniformly
in \(j,T,M\). For a fixed target \(A>2\), write
\(K_{\mathrm{SW}}(A)\) for the
chosen logarithmic level. For auxiliary moduli \(k\le L^{K_{\mathrm{SW}}(A)}\), this is prime
Siegel--Walfisz at arbitrarily high fixed logarithmic precision.
The divisor loss from excluding primes dividing the auxiliary integers is
absorbed into \(d(h)^{B_{\mathrm{div}}}\). For \(k>L^{K_{\mathrm{SW}}(A)}\), use
\[
 |\Delta|\le |J|/k+1+\pi(J)/\varphi(k)
\]
and \(\sup_{k>L^{K_{\mathrm{SW}}(A)}}\varphi(k)^{-1}
\ll L^{-K_{\mathrm{SW}}(A)}\log\log(3L^{K_{\mathrm{SW}}(A)})\), taking
\(K_{\mathrm{SW}}(A)\) larger than the
requested saving. For every fixed target \(A>2\), choose
\(H_g=A+4\). The grid has \(O_b(L^{H_g+1})\) intervals. The prime
number theorem gives the norm bounds \eqref{eq:k2-interval-norms}.
With \(A_{\mathrm{block}}=A+H_g/2+2\), Fact~\ref{fact:bv-full-support}
and Paper~II, (5.23)--(5.28), bound the rectangular blocks by
\(O_{A,b}(NL^{-A-3/2})\). The boundary term in
\eqref{eq:k2-boundary-cost} is \(O(NL^{-A-3}+Q)\). Increasing the
single fixed logarithmic level buffer absorbs the \(x\asymp_b N\)
constants. Thus the complete moving-top source, including the
actual top \(T=N'\), satisfies \eqref{eq:k2-bv-source}. No fixed-width
truncation estimate is substituted.


\subsection{Four directed sieve rows}

Let \(\mathcal P_N=\{p\in\mathbb P:p\nmid2N\}\).
Put
\begin{equation}\label{eq:k2-sifted-source}
 S_b(d;z)=\sum_{d\mid N-e}A_b(e)\,1_{((N-e)/d,\,P_N(z))=1},\qquad
 P_N(z)=\prod_{p<z,\ p\in\mathcal P_N}p .
\end{equation}
There is no prefix--cofactor coprimality condition.
Both source primes exceed \(N^{\mu_*}\) eventually.
Thus, for every supported modulus \(d\),
\begin{equation}\label{eq:k2-local-density}
 \sum_{d\mid N-e}A_b(e)
 =\frac{X_b}{\varphi(d)}+\Delta_{A_b}(d,N),\qquad
 (d,2N)=1,\quad P^+(d)<N^{\mu_*}.
\end{equation}
Here \(P^+(d)\) is the largest prime factor, with \(P^+(1)=1\).
For the squarefree sieve divisors, use the multiplicative density
\begin{equation}\label{eq:k2-squarefree-density}
 g(p)=\frac1{p-1}\quad(p\in\mathcal P_N),\qquad
 g(p^j)=0\quad(j\ge2),\qquad
 g(p)=0\quad(p\notin\mathcal P_N).
\end{equation}
This is the squarefree-divisor density of the linear-sieve application.
That input requires its zero prime-power tail. The
full arithmetic density is \(g(p^j)=1/\varphi(p^j)\) for
\(p\in\mathcal P_N\) and \(j\ge1\). No estimate uses its prime-power
tail. The two densities agree at active primes. These are the only local
factors used by the main-mass calculation.
In Paper~I, Theorem~2.6 notation, this density is
\[
 \omega(p)=\frac p{p-1}\quad(p\in\mathcal P_N),\qquad
 \omega(p)=0\quad(p\notin\mathcal P_N),\qquad
 \omega(p^a)=0\quad(a\ge2).
\]
Thus \(\omega(p)/p=g(p)\), and condition (2) has zero left side. For
\(p\ge3\),
\[
 \left(1-\frac1{p-1}\right)^{-1}
 =\left(1-\frac1p\right)^{-1}
   \left(1+\frac1{p(p-2)}\right).
\]
The product of the second factors over \(p\ge w\) is
\(1+O(1/w)=1+O(1/\log w)\). Hence Mertens' product estimate gives
condition (1) with exponent one and one fixed \(K_{\mathrm{dim}}\ge2\).
We may take \(L_{\omega}=2\). These are the dimension and prime-power
inputs of Iwaniec's linear sieve \cite{Iwaniec1980Remainder}.

If \(N-e\) is prime, it is at least \(NL^{-333}\).
It therefore exceeds \(N^{\mu_*}\) eventually.
Hence
\begin{equation}\label{eq:k2-source-comparison}
 B_b'\le S_b(1;N^{\mu_*}).
\end{equation}

Use the ordered primes
\(N^u\le a_j<\cdots<a_1<N^{\mu_*}\)
in \(\mathcal P_N\).
The exact decomposition of
\eqref{eq:source-buchstab}--\eqref{eq:source-upper-slack} gives
\begin{equation}\label{eq:k2-buchstab-upper}
 B_b\le E_{\mathrm{top}}+T_0-T_1+T_2-R_3+P_3.
\end{equation}
Here
\begin{equation}\label{eq:k2-four-rows}
\begin{aligned}
 T_0&=S_b(1;N^u),\\
 T_j&=\sum_{N^u\le a_j<\cdots<a_1<N^{\mu_*}}
                 S_b(a_1\cdots a_j;N^u)\quad(j=1,2),\\
 R_3&=\sum_{\mathrm{ret}}S_b(a_1a_2a_3;N^{t_3}),\\
 P_3&=\sum_{\mathrm{ret}}\sum_{\substack{N^{t_3}\le r<a_3\\r\in\mathcal P_N}}
                       S_b(a_1a_2a_3r;N^{t_3}).
\end{aligned}
\end{equation}
Every prime in these sums belongs to \(\mathcal P_N\).
The last sum keeps the half-open cells it inherits.
The symbol \(\sum_{\mathrm{ret}}\) restricts by
\begin{equation}\label{eq:k2-retention}
 \Theta_{\mathfrak c}=\inf_{\overline{\mathfrak c}}H,\qquad
 \ell_{\mathfrak c}=\Theta_{\mathfrak c}-\eta_L-t_1-t_2-t_3,
 \qquad \ell_{\mathfrak c}\ge2\rho t_3.
\end{equation}
Here \(H\) is the full eight-branch packing function,
\eqref{eq:classical-packing-ceiling}.
The infimum uses the complete closed ordered cell.
It does not use a quadrature node.

The four directed applications are
\begin{equation}\label{eq:k2-aa-directions}
\begin{array}{c|c|c|c}
 \text{term}&\text{row sign}&\text{sieve direction}&\text{level exponent}\\ \hline
 T_0&+&\text{upper }F&1/2-\eta_L\\
 T_1&-&\text{lower }f&1/2-\eta_L-\rho t_1\\
 T_2&+&\text{upper }F&1/2-\eta_L-\rho(t_1+t_2)\\
 R_3&-&\text{lower }f&\ell_{\mathfrak c}
\end{array}
\end{equation}
The first three cutoffs are \(N^u\).
The fourth cutoff is \(N^{t_3}\).
The canonical functions are Iwaniec's linear-sieve functions
\cite{Iwaniec1980Remainder}.
Extend \(f\) by zero for \(s\le2\).
No shifted beta-sieve profile is substituted.

For each prefix product \(h\), let \(\mathcal A_h\) be the cofactor
sequence indexed by \(n=(N-e)/h\), with multiplicity \(A_b(e)\) when
\(h\mid N-e\). Put \(X_h=X_b/\varphi(h)\). Write \(z_h\) for the
row cutoff, \(D_h\) for its level, and \(s_h=\log D_h/\log z_h\).
Every prefix prime is at least \(z_h\), and each squarefree sieve
modulus \(t\mid P_N(z_h)\) has prime factors below \(z_h\). Thus
\((h,t)=1\), and the local identity \eqref{eq:k2-local-density} gives
the exact remainder identity
\begin{equation}\label{eq:k2-aa-prefix-remainder}
 r(\mathcal A_h,t)=\Delta_{A_b}(ht,N).
\end{equation}
Since \(\omega(t)/t=1/\varphi(t)\) and \((h,t)=1\), the local main
term is \(X_h\omega(t)/t=X_b/\varphi(ht)\), as required by
Theorem~2.6. For rows \(j=0,1,2\), put
\(S=t_1+\cdots+t_j\), with \(S=0\) for row zero, and
\(\theta=1/2-\eta_L\). The clipped exponent cells give
\[
 h\le N^{\rho S},\qquad D_h=N^{\theta-\rho S},\qquad
 S\le2\mu_*.
\]
Hence \(\log D_h/\log N\ge\theta-2\rho\mu_*>0.1092>2u\).
The upper rows satisfy \(N^u\le D_h\). The lower row satisfies
\(N^u\le D_h^{1/2}\). For sufficiently large
\(N\), with \(\sigma\) fixed, \(D_h\ge2\).
Apply Paper~I, Theorem~2.6, (2.30), with density
\(g(p)=1/(p-1)\) on \(\mathcal P_N\), \(g(p^a)=0\) for \(a\ge2\),
and sieve parameter \(\epsilon_s=\sigma\). With
\(\overline E=2C_{\mathrm{AA37}}\sigma\), the resulting directed
inequalities are
\begin{equation}\label{eq:k2-aa-directed-estimates}
\begin{aligned}
 S(\mathcal A_h,z_h)
 &\le X_hV_N(z_h)\{F(s_h)+E^+_{s_h}\}+R_1^+(h)+R^+(h),\\
 S(\mathcal A_h,z_h)
 &\ge X_hV_N(z_h)\{f(s_h)-E^-_{s_h}\}-R_1^-(h)-R^-(h),\\
 |E^{\pm}_{s_h}|&\le\overline E.
\end{aligned}
\end{equation}
Only the upper bound \(|E^{\pm}_{s_h}|\le\overline E\) is used. No sign of
\(E^{\pm}_{s_h}\) is assumed. Since \(X_hV_N(z_h)\ge0\), replacing
\(E^{\pm}_{s_h}\) by \(\overline E\) weakens each inequality.
The actual dimension and prime-power constants, the level lower bound,
and the fixed-\(\sigma\) threshold are in
\eqref{eq:k2-aa-decay}.

The superscript \(\mathrm{AA}\) below marks errors from this
linear-sieve theorem. It does not denote a separate family of profiles.
For row \(j\), sum over its actual prefix set \(\mathcal H_j\). Let
\(\mathcal F_j\) be \(F\) on the upper rows and \(f\) on the lower
rows. The definitions used in the signed equalities are
\begin{equation}\label{eq:k2-aa-row-definitions}
\begin{aligned}
 m_j&=\sum_{h\in\mathcal H_j}X_hV_N(z_h)\mathcal F_j(s_h),\\
 \mathfrak M_j&=\sum_{h\in\mathcal H_j}X_hV_N(z_h),\\
 \mathcal A_j^{\mathrm{AA}}&=\overline E\,\mathfrak M_j,\\
 r_j&=\sum_{h\in\mathcal H_j}\{R_1^{\nu_j}(h)+R^{\nu_j}(h)\},
 \qquad \nu_j=+\text{ on upper rows},\quad \nu_j=-\text{ on lower rows},\\
 z_j&=m_j+\mathcal A_j^{\mathrm{AA}}+r_j-T_j
       \quad\text{on upper rows},\\
 z_j&=T_j-m_j+\mathcal A_j^{\mathrm{AA}}+r_j
       \quad\text{on lower row }j=1,\\
 z_3&=R_3-m_3+\mathcal A_3^{\mathrm{AA}}+r_3.
\end{aligned}
\end{equation}
Summing \eqref{eq:k2-aa-directed-estimates} over prefixes proves
\(z_j\ge0\). Also \(\mathfrak M_j\ge0\), so
\(\mathcal A_j^{\mathrm{AA}}\ge0\). These nonnegativity
statements are proved at source level for the exact signed remainders in
\eqref{eq:k2-aa-prefix-remainder}. No sign of these remainders is assumed.
Write the directed equalities as
\begin{equation}\label{eq:k2-aa-equalities}
\begin{aligned}
 T_0&=m_0+\mathcal A_0^{\mathrm{AA}}+r_0-z_0,&
 T_1&=m_1-\mathcal A_1^{\mathrm{AA}}-r_1+z_1,\\
 T_2&=m_2+\mathcal A_2^{\mathrm{AA}}+r_2-z_2,&
 R_3&=m_3-\mathcal A_3^{\mathrm{AA}}-r_3+z_3,
 \qquad \mathcal A_j^{\mathrm{AA}},z_j\ge0.
\end{aligned}
\end{equation}
\begin{lemma}[Signed row identity]\label{lem:k2-row-identity}
With \(M_K=m_0-m_1+m_2-m_3\),
\[
 \Delta_{\mathrm{arith}}:=T_0-T_1+T_2-R_3-M_K
 =\sum_{j=0}^3\mathcal A_j^{\mathrm{AA}}+\sum_{j=0}^3r_j-\sum_{j=0}^3z_j,
\]
and every \(\mathcal A_j^{\mathrm{AA}}\) occurs with coefficient \(+1\).
\end{lemma}
\begin{proof}
Substitute \eqref{eq:k2-aa-equalities}. The lower rows carry the row
sign \(-\), which turns \(-\mathcal A_j^{\mathrm{AA}}-r_j+z_j\) into
\(+\mathcal A_j^{\mathrm{AA}}+r_j-z_j\).
\end{proof}
For rows \(j=0,1,2\), the total moduli in these remainders fit the
level of \eqref{eq:k2-bv-source}. A tail tuple satisfies
\(\prod_iD_i<D_h\). In the parity ending in a defining cubic condition,
remove two copies of the last label. In the other parity, use the
preceding cubic condition and the ordering of the last two labels. The
empty and one-label cases follow from the tuple definitions.
Each \(p_i\) lies in the interval
\(D_i\le p_i<\min(D_i^\rho,z_h)\), and
\(d<D_h^\sigma\). Thus, with the same \(S,\theta\) as above,
\[
 q=h d\prod_i p_i\le hD_h^{\sigma+\rho}
 \le N^{\rho S+(\sigma+\rho)(\theta-\rho S)}
 =N^{\theta+(\sigma+\tau)(\theta-\rho S)}
 \le N^{\theta+(\sigma+\tau)\theta}
 <N^{1/2-\eta_L/2}.
\]
This power buffer is below \(N^{1/2}L^{-B}\) for every fixed \(B\),
once \(N\) is large enough.

The factors of such a modulus that are at least \(z_h\) are exactly the
ordered prefix primes. They determine \(h\) and its clipped cell.
Factoring the remaining part at \(D_h^{\sigma^2}\) determines the
squarefree smooth divisor and the tail-prime multiset. The number of
tail labels is at most \(J_\sigma=\lceil\sigma^{-2}\rceil\).
Equal-label permutations contribute at most \(J_\sigma!\). Each
AA coefficient has modulus at most one. The number of label vectors
is less than \(\exp(8\sigma^{-3})\). Hence the total coefficient at each
modulus is at most
\(C_{\mathrm{K2},\sigma}=1+\exp(8\sigma^{-3})J_\sigma!\), independent of \(N\).
Apply \eqref{eq:k2-bv-source} with target \(A=4\). Divide by
\(\mathcal N_K=(1-\varepsilon)C(N)N/L^2\); here \(C(N)\ge C_2>0\)
by Paper~II, (1.2). This proves
\begin{equation}\label{eq:k2-low-row-remainders}
 \frac{|r_0|+|r_1|+|r_2|}{\mathcal N_K}\ll_{b,\sigma}L^{-2}.
\end{equation}
The low-row bound uses Fact~\ref{fact:bv-full-support}, the
prime Siegel--Walfisz theorem, the prime number theorem and the Paper~II
convolution inputs. It gives no estimate for the principal
masses or row three.

The \(m_j\) are the directed main terms.
The \(r_j\) are signed arithmetic remainders.
Thus, with \(M_K=m_0-m_1+m_2-m_3\),
\begin{equation}\label{eq:k2-signed-join}
\begin{aligned}
 B_b\le{}&E_{\mathrm{top}}+P_3+M_K
       +\sum_{j=0}^3\mathcal A_j^{\mathrm{AA}}+\sum_{j=0}^3r_j-\sum_{j=0}^3z_j\\
 \le{}&E_{\mathrm{top}}+P_3+M_K
       +\sum_{j=0}^3\mathcal A_j^{\mathrm{AA}}+\sum_{j=0}^3|r_j|.
\end{aligned}
\end{equation}
This step is finite algebra.
The directed source comparisons are proved separately.

The canonical row masses are explicit.
For a prefix cell \(\mathfrak c\), put
\[
 H_j(\mathfrak c)=
 \sum_{\substack{a_1>\cdots>a_j\\a_i\text{ in the cell}}}
       \frac1{\varphi(a_1\cdots a_j)}.
\]
Since the prefix primes exceed every sieve prime, the product over sieve
primes at each row is \(V_N(z_h)\). With the ordered cell sets
\(D_j\) from
\eqref{eq:k2-parent-functional},
\begin{equation}\label{eq:k2-row-main-masses}
\begin{aligned}
 m_0&=X_bV_N(N^u)F(s_0),\\
 m_1&=X_bV_N(N^u)\sum_{\mathfrak c\in D_1}
                 H_1(\mathfrak c)f(s_1(\mathfrak c)),\\
 m_2&=X_bV_N(N^u)\sum_{\mathfrak c\in D_2}
                 H_2(\mathfrak c)F(s_2(\mathfrak c)),\\
 m_3&=X_b\sum_{\mathfrak c\in D_3}\chi_{\mathfrak c}
       V_N(N^{t_3})H_3(\mathfrak c)f(\ell_{\mathfrak c}/t_3).
\end{aligned}
\end{equation}
The masses \(\mathfrak M_j\) are these sums with \(F\) and \(f\)
replaced by \(1\). In particular, the row-three retention indicator
\(\chi_{\mathfrak c}\) appears in both sums.


\subsection{The principal functional}

We display the numerical functional used for \(M_K\).
Set
\begin{equation}\label{eq:k2-canonical-profile}
\begin{aligned}
 v_1(s)&=\min\left\{\frac{1+s}{2},
             \frac{32-34s}{39},\frac{57}{64}-\frac32s\right\},\\
 v(s)&=v_1(s)-u,\qquad \psi(x,y)=x1_{x\ge y},\\
 \Theta_*(x,y,z)&=\max\{v(x),v(y),v(x+y),v(x+y+z),\\
 &\hspace{28mm}\psi(v(x+z),x+2y+z),
                   \psi(v(y+z),2x+y+z)\}.
\end{aligned}
\end{equation}
The subtraction of \(u\) precedes both guards.
This is Paper~II, (4.70)--(4.72)
\cite[(4.70)--(4.72)]{PaperII}.
Put
\begin{equation}\label{eq:k2-g-functional}
\begin{aligned}
 H_{\log}&=\log(\mu_*/u),\\
 R_{\mathrm{up}}&=500(F(30)-1)+500H_{\log}(1-f(30))
                    +250H_{\log}^2(F(30)-1),\\
 J_*&=\int_u^{\mu_*}\frac{dx}{x}
         \int_u^x\frac{dy}{y}
         \int_u^y\frac{1-f((\Theta_*(x,y,z)-x-y-z)/z)}{z^2}\,dz,\\
 G_*&=\frac1{\mu_*}+R_{\mathrm{up}}+J_*.
\end{aligned}
\end{equation}
We keep the true delay-function tail.
Paper~II, (4.73)--(4.75), gives the constant cancellation
\begin{equation}\label{eq:k2-constant-cancellation}
 \frac{1-H_{\log}+H_{\log}^2/2}{u}
 -\int_u^{\mu_*}\frac{dx}{x}\int_u^x\frac{dy}{y}
                         \int_u^y\frac{dz}{z^2}
 =\frac1{\mu_*}.
\end{equation}

The finite-parent comparison also needs its direction.
Fix a finite clipped partition \(\Pi\) of \([u,\mu_*)\).
Its cells have endpoint ratios at most \(\rho\).
Each clipped cell is a product of half-open exponent intervals \([t_i,t_i')\) with \(t_i'\le\rho t_i\). Here \(t_i\) is the lower endpoint, so \(N^{t_i}\le a_i<N^{\rho t_i}\), i.e. \(t_i\le x_i<\rho t_i\). The closure \(\overline{\mathfrak c}\) is used only for infima and lower bounds. The cutoffs \(N^{t_3}\) in \(R_3,P_3\) and the factors \(\rho t_1,\rho(t_1+t_2)\) in \(s_1,s_2\) use this convention.
For a clipped cell \(\mathfrak c\), let
\(w_{\mathfrak c}=\int_{\mathfrak c}\prod_i dx_i/x_i\).
Define
\begin{equation}\label{eq:k2-parent-functional}
\begin{aligned}
 s_0&=(1/2-\eta_L)/u,\\
 s_1(\mathfrak c)&=(1/2-\eta_L-\rho t_1)/u,\\
 s_2(\mathfrak c)&=(1/2-\eta_L-\rho(t_1+t_2))/u,\\
 E_\Pi&=\frac{F(s_0)-1}{u}
 +\sum_{\mathfrak c\in D_1}\frac{w_{\mathfrak c}}u(1-f(s_1(\mathfrak c)))\\
 &\qquad+\sum_{\mathfrak c\in D_2}\frac{w_{\mathfrak c}}u(F(s_2(\mathfrak c))-1),\\
 P_\Pi&=\frac{1-H_{\log}+H_{\log}^2/2}{u}+E_\Pi
 -\sum_{\mathfrak c\in D_3}
    w_{\mathfrak c}1_{\ell_{\mathfrak c}\ge2\rho t_3}
             \frac{f(\ell_{\mathfrak c}/t_3)}{t_3}.
\end{aligned}
\end{equation}
Here \(D_j=\{u\le x_j<\cdots<x_1<\mu_*\}\).
Each cell is intersected with this ordered domain.
Its weight includes the actual ordered multiplicity.
The required comparisons are
\begin{equation}\label{eq:k2-parent-comparison}
 0\le E_\Pi\le R_{\mathrm{up}},\qquad
 P_\Pi\le\frac1{\mu_*}+E_\Pi+J_*\le G_*.
\end{equation}
The last line is an upper comparison.
It asserts no equality between finite sums and integrals.

Compare the six canonical branches of \(\Theta_*\) with their matching
branches among the eight parent branches of the packing function \(H\) in
\eqref{eq:classical-packing-ceiling}. Adding \(u\) to an unguarded canonical
branch gives the matching parent branch. For an active canonical guarded
branch, the parent guard is also active, and adding \(u\) again gives the
matching parent branch. An inactive canonical guarded branch is zero and
cannot maximize the canonical ceiling, since its first unguarded branch is
positive on the ordered domain. The two extra parent branches can only
increase \(H\). Hence
\(H(x,y,z)\ge\Theta_*(x,y,z)+u\) pointwise on that domain. On a closed
anchored cell, each coordinate moves by at most \(\mu_*\tau\). Fix a point
and a maximizing canonical branch there, and use its matching parent branch.
If it is guarded, its guard has reserve \(u-7\mu_*\tau>0\) throughout the
cell. This selected branch changes by at most \((9/2)\mu_*\tau\) across the
cell. Therefore
\begin{equation}\label{eq:k2-parent-reserve}
\begin{aligned}
 \Theta_{\mathfrak c}&\ge\Theta_*(x,y,z)+u-(9/2)\mu_*\tau,\\
 \ell_{\mathfrak c}&\ge\Theta_*(x,y,z)-x-y-z+g,\\
 g&=u-\eta_L-(9/2)\mu_*\tau>2\mu_*\tau.
\end{aligned}
\end{equation}
Monotonicity of \(f(s)\), together with this reserve,
gives the directed kernel comparison
\begin{equation}\label{eq:k2-parent-kernel}
 \chi_{\mathfrak c}\frac{f(\ell_{\mathfrak c}/t_3)}{t_3}
 \ge
 \frac{f((\Theta_*(x,y,z)-x-y-z)/z)}{z}.
\end{equation}
Summing over cells gives the upper comparison in
\eqref{eq:k2-parent-comparison}. It does not identify a finite sum with
the integral.


Put \([x]_+=\max\{x,0\}\).
We need the arithmetic-to-integral estimate
\begin{equation}\label{eq:k2-main-mass}
 \frac{[M_K-\mathcal N_K\,2e^{-\gamma}I(b)G_*]_+}
      {\mathcal N_K}=o_{b,\sigma}(1).
\end{equation}
The quantity \(M_K\) contains only the canonical values of \(F\) and
\(f\). The linear-sieve errors \(E^{\pm}_{s_h}\) and the remainders
\(r_j\) lie outside it. They are paid in
\eqref{eq:k2-aa-mass} and \eqref{eq:k2-remainder-estimates}. Fix
\((b,\sigma)\). Three inputs enter. The first is the two-orientation
arithmetic-mass estimate for the source \eqref{eq:k2-source}. The second
is its four signed canonical AA
main masses \eqref{eq:k2-row-main-masses}. The third is the finite functional
\(P_\Pi\) in \eqref{eq:k2-parent-functional}. For all sufficiently
large even \(N\) they give
\[
 M_K/\mathcal N_K\le2e^{-\gamma}I(b)G_*+o_{b,\sigma}(1).
\]
Here the error tends to zero with \(b,\sigma\) and the cell count fixed,
as \(N\) tends to infinity. This estimate is for the finite functional in
\eqref{eq:k2-parent-functional}. It gives
\eqref{eq:k2-main-mass} directly, so
\(N_{\mathrm{mass}}(b,\sigma)\) can be the threshold of
this fixed-\((b,\sigma)\) statement. This bound controls only the
principal-mass term. It does not control the AA remainders, the optional
numerical certificate, the complete exterior comparison or
the endpoint terms. No sharper arithmetic-mass estimate is used here.


For the exterior comparison we use the common $G_0$ coefficient
\begin{equation}\label{eq:k2-certificate}
\begin{gathered}
 G_0=\frac{6088440363}{10^9},\qquad
 G_*\le G_0,\\
 2e^{-\gamma}G_0\le c_{K2,G0}
 :=\frac{6836825163875517756286846321650754372931094689}{10^{45}}.
\end{gathered}
\end{equation}
For \(G_*\le G_0\), use the tight endpoint
\(J_*\le0.967120\) of Paper~II, (A.10a). It holds for the exact profile
\eqref{eq:k2-g-functional} \cite[(4.73), (4.75), (A.10a)]{PaperII}.
Its weaker displayed endpoint is not used.
The Taylor-model delay certificate is used for this coefficient comparison
and is not an input to the shifted replay itself. It gives the explicit bounds
\[
 0\le F(30)-1<3\cdot10^{-49},\qquad
 0\le1-f(30)<3\cdot10^{-49}.
\]
Since \(0<H_{\log}=\log(3125/32)<5\),
\eqref{eq:k2-g-functional} then gives
\[
 R_{\mathrm{up}}< (500+500\cdot5+250\cdot25)3\cdot10^{-49}
 <10^{-44},\qquad
 G_*<\frac{128}{25}+\frac{967120}{10^6}+10^{-44}
 <\frac{6088440363}{10^9}=G_0.
\]
The finite delay calculation and this rational budget do not replace the
source-to-profile comparison or its fixed-parameter hypotheses.
For the scalar inequality, define
\begin{equation}\label{eq:k2-gamma-enclosure}
\begin{gathered}
 \gamma_-:=\frac{57721566490153286060651209008240243104215933593992}{10^{50}},\\
 \gamma_+:=\frac{57721566490153286060651209008240243104215933593993}{10^{50}},\\
 \gamma_-<\gamma<\gamma_+,\\
 c_{K2,G0}-2G_0\sum_{j=0}^{60}\frac{(-\gamma_-)^j}{j!}
       >\frac{75307}{10^{50}}>0.
\end{gathered}
\end{equation}
The rational enclosure follows from Euler--Maclaurin for \(H_n\) at
\(n=10^5\), through \(B_{16}\). The next-term remainder is bounded
by \(|B_{18}|/(18n^{18})\). Both \(\log n\) and \(H_n\) are enclosed
outwards by the certified exact-input record. For \(0<x<1\), the even Taylor
partial sum of \(e^{-x}\) is an upper bound. Thus the exact rational
comparison in \eqref{eq:k2-gamma-enclosure} proves
\(2e^{-\gamma}G_0<c_{K2,G0}\).


The switched comparison uses the same coefficient.
At \(b=2999993/7000000\), the common final-margin inequality
\eqref{eq:bill-margin-simple} has slack
\(0.000228954311267289\ldots>1/5000\) on the stated route \(g_6:=0\).
Both comparisons use the same $G_0$ bound. Scalar positivity does not
establish the required source estimates.

\subsection{The complete error bound}

The analytic hypotheses are stated below.
The sums are over the original applications.
The four principal linear-sieve errors satisfy
\begin{equation}\label{eq:k2-aa-mass}
\begin{gathered}
 \frac{\sum_{j=0}^3 \mathcal A_j^{\mathrm{AA}}}{\mathcal N_K}
 \le A_K(b)C_{\mathrm{AA}}\sigma,\\
 A_K(b)=2e^{-\gamma}I(b)W_{\mathrm{AA}},\qquad
 W_{\mathrm{AA}}=500\left(1+H_{\log}
                    +\frac{H_{\log}^2}2+\frac{H_{\log}^3}6\right).
\end{gathered}
\end{equation}
Let \(C_{\mathrm{AA37}}\) be the absolute constant in Paper~I,
Theorem~2.6, (2.31), and define
\begin{equation}\label{eq:k2-aa-constant}
 C_{\mathrm{AA}}=4C_{\mathrm{AA37}}.
\end{equation}
The coefficient four is the factor two in
\eqref{eq:k2-aa-density-mass} times the factor two in the common bound
\(|E_s^{\pm}|\le 2C_{\mathrm{AA37}}\sigma\). The sum over the four rows is
already included in \(W_{\mathrm{AA}}\). It is independent of \(\sigma,N\).
It is unrelated to the global tensor constant in
\eqref{eq:ra-global-tensor}. All four actual levels satisfy
\(D\ge N^{1/250}\). We use (2.31) with the actual fixed dimension constant
\(K_{\mathrm{dim}}\). For the squarefree density in
\eqref{eq:k2-squarefree-density}, the left side of (2.27) is zero. Hence
the actual prime-power constant is \(L_{\omega}=2\). This value is used
in the exponential below.
\begin{equation}\label{eq:k2-aa-decay}
 |E_s^{\pm}|\le C_{\mathrm{AA37}}
 \{\sigma+\sigma^{-8}e^{K_{\mathrm{dim}}+2}(\log D)^{-1/3}\}
 \le2C_{\mathrm{AA37}}\sigma
\end{equation}
once \(L\ge250\exp(3(K_{\mathrm{dim}}+2))\sigma^{-27}\).
Name the corresponding even threshold
\[
 N_{\mathrm{AA,dec}}(\sigma)
 :=2\left\lceil\frac12
 \exp\!\left(250\exp(3(K_{\mathrm{dim}}+2))\sigma^{-27}\right)
 \right\rceil.
\]
It is chosen at fixed \(\sigma\). Every
\(N\ge N_{\mathrm{AA,dec}}(\sigma)\) satisfies the displayed decay bound.
The four ordered harmonic row masses, before the common Euler and source
factor, are bounded by
\begin{equation}\label{eq:k2-aa-row-masses}
\begin{gathered}
 \frac1u,\qquad \frac{H_{\log}}u,\qquad
 \frac{H_{\log}^2}{2u},\\
 \text{row-3 mass}\le
 \rho\int_u^{\mu_*}\frac{dx}{x}\int_u^x\frac{dy}{y}
        \int_u^y\frac{dz}{z^2}
 \le\frac{\rho H_{\log}^2}{2u}
 <\frac{H_{\log}^3}{6u}.
\end{gathered}
\end{equation}
For the last inequality use \(3\rho<H_{\log}=\log(3125/32)\).
The factor \(\rho\) allows the anchored-bin endpoint enlargement. The first three
identities are the volumes of ordered logarithmic simplexes. For the
fourth, bound \(\int_u^y dz/z^2\le1/u\) and integrate over the ordered
\((x,y)\) simplex. The sum of the four displayed upper envelopes is \(W_{\mathrm{AA}}\),
the four-row envelope used here. It is distinct from Paper~II's
coarser global sieve-error bound.
Let \(H=H_{\log}=\log(3125/32)\). PNT and Mertens' theorems, uniformly
over the finitely many fixed-\(\sigma\) cells, give
\[
 V_N(N^a)=\frac{2e^{-\gamma}C(N)}{aL}(1+o(1)),\qquad
 H_j(\mathfrak c)=w_{\mathfrak c}+o(1).
\]
For the retained row, \(1/t_3\le\rho/x_3\) on each cell. Dropping
the retention restriction only increases its mass. Consequently
\begin{equation}\label{eq:k2-aa-density-derived}
\begin{aligned}
 \frac{\sum_{j=0}^3\mathfrak M_j}{\mathcal N_K}
 &\le2e^{-\gamma}I(b)
 \left\{\frac1u+\frac{H}{u}+\frac{H^2}{2u}+\rho J_4\right\}+o(1),\\
 J_4&=\frac{H^2/2-H+1}{u}-\frac1{\mu_*}
       \le\frac{H^2}{2u}.
\end{aligned}
\end{equation}
The bracket is smaller than
\(W_{\mathrm{AA}}=u^{-1}(1+H+H^2/2+H^3/6)\), with relative gap greater
than \(0.2842\) for \(0<\sigma\le10^{-10}\). The coarser bound
\(\rho J_4\le\rho H^2/(2u)\) still leaves relative gap greater than \(0.1723\).
Consequently, for each fixed \(b,\sigma\), there is a threshold
\(N_{\mathrm{AA,dens}}(b,\sigma)\) such that, for every even
\(N\ge N_{\mathrm{AA,dens}}(b,\sigma)\),
\(\sum_j\mathfrak M_j\le\mathcal N_K\,2e^{-\gamma}I(b)W_{\mathrm{AA}}\).
Define
\[
 N_{\mathrm{AA}}(b,\sigma)
 :=\max\{N_{\mathrm{AA,dens}}(b,\sigma),N_{\mathrm{AA,dec}}(\sigma)\}.
\]
Thus both the density estimate and \eqref{eq:k2-aa-decay} hold above
\(N_{\mathrm{AA}}(b,\sigma)\), as required for \eqref{eq:k2-aa-mass}.
The maxima defining \(T_K\) and \(T_r\) below include this combined threshold.
The displayed source-density bound below is weaker by a factor of two.
Let \(\mathfrak M_j\) be the nonnegative principal mass multiplying the
relative sieve error in row \(j\). The resulting bound is
\begin{equation}\label{eq:k2-aa-density-mass}
 \sum_{j=0}^3\mathfrak M_j\le
 2\mathcal N_K\,2e^{-\gamma}I(b)W_{\mathrm{AA}}.
\end{equation}
Multiplying it by
\(2C_{\mathrm{AA37}}\sigma\) gives \eqref{eq:k2-aa-mass} with
\eqref{eq:k2-aa-constant}. The Euler factors and the actual-source mass
identification are included above. The secondary decay is paid once.


The positive peel is bounded by the two-form argument of Paper~II,
\cite[(4.46)--(4.56)]{PaperII}, applied to the original source.
Put \(\mathcal B_m=\{e:\ m\mid N-e\}\) with multiplicity \(A_b(e)\).
For \(m=a_1a_2a_3r\) in \(P_3\), put \(E=N^{t_3}\) and \(w=N^u\le E\).
Write \(A_b=\sum_{h=1}^21_{\{A_b\ge h\}}\).
Every \(e\) in the support is a product of two primes at least \(N^b\).
Hence \(e\) is odd and \(w\)-rough. Write \(N-e=mk\). Then \(k\) is odd. A
prime \(p<w\) dividing \(k\) and \(N\) would divide \(e\), which is impossible.
Hence both forms \(k\) and \(N-mk\) are free of all primes below \(w\).
Therefore \(S_b(m;E)\) is at most the two-form Selberg majorant of
\cite[(4.48)--(4.53)]{PaperII} on \(1\le k\le N/m\), once for each layer.
All prime factors of \(m\) are at least \(E\ge w\) and lie in
\(\mathcal P_N\). Hence \((m,N)=1\), and \(\nu_N(p)\) is as in
\cite[(4.48)]{PaperII}. Retention gives, as in \cite[(4.55)]{PaperII},
\begin{equation}\label{eq:k2-peel-ranges}
 \log_Nm\le\rho(t_1+t_2+2t_3)\le\rho(\Theta_{\mathfrak c}-\eta_L)
 <\tfrac{153}{256}+2\tau,\qquad
 \log_N(w^2m)<1-0.0116,\qquad \log_Nw=u\ge\tfrac2{1001}.
\end{equation}
Consequently \cite[(4.47), (4.53), (4.54)]{PaperII} give, eventually and
uniformly in \(m\),
\begin{equation}\label{eq:k2-peel-inputs}
 S_b(m;E)\le2\cdot\frac{4.002}{u^2}\cdot\frac{U_C}{m}
 \le\frac{2008008\,U_C}{m},\qquad
 \max_hM_h\le2\tau,\qquad \sum_hM_h\le\frac{459}{100},
\end{equation}
since one layer is at most \((N/m)/\mathcal G_N(w)+O(w^2\log^5w)\) with
\(\mathcal G_N(w)=(1+o(1))u^2L^2/(4C(N))\), and
\(4.002/u^2=1000500<1004004=2\cdot502002\). The bin sums are
\cite[(4.56)]{PaperII}: \(\log(\mu_*/u)=\log(3125/32)<4.59\), and
\(a_3,r\) lie in one parent bin. Therefore
\begin{equation}\label{eq:k2-peel}
 \frac{P_3}{U_C}\le2008008\cdot2\left(\frac{459}{100}\right)^3\tau
 =\frac{24272444031579}{62500}\tau<388400000\,\sigma^9,
\end{equation}
and \(P_3/\mathcal N_K\le P_K\sigma^9\) with
\(P_K=388400000/(1-\varepsilon)=3884000000000/9999\), for
\(N\ge N_{\rm peel}(\sigma)\). The coefficient \(2008008\) includes the two
layers. No orientation factor is added.

The remaining estimates are
\begin{equation}\label{eq:k2-remainder-estimates}
\begin{aligned}
 \frac{|r_0|+|r_1|+|r_2|}{\mathcal N_K}&\ll_{b,\sigma}L^{-2},\\
 \frac{|r_3|}{\mathcal N_K}&\ll_{b,\sigma}
 N^{-\ell_0\sigma^2}L^3+N^{-b/4}L^7+L^{-260}+L^{-120}
 +L^{-2}+N^{-c_\sigma}L^{4024+C}
 =O_{b,\sigma}(L^{-2}).
\end{aligned}
\end{equation}
The first line is \eqref{eq:k2-low-row-remainders}. The second is
Proposition~\ref{prop:k2-row3} below. Its displayed terms come, in order,
from the squareful mask, proper powers, logarithm freezing and the
boundary. The last two are the combined logarithmic Type-II and Type-I-2
bound and the Type-II power-saving term.


\subsection{The row-three remainder}\label{subsec:k2-row3}

We bound the complete third-row remainder of the actual source
\eqref{eq:k2-source} by Paper~I, Theorem~3.2. Every other contribution is
bounded separately. Fix \(b,\sigma\). Write \(\Xi^{\sharp}_3\) for the
retained cells \(\mathfrak c\in D_3\) with \(\chi_{\mathfrak c}=1\). Its
size is \(O_\sigma(1)\). For \(\mathfrak c\in\Xi^{\sharp}_3\) put
\(D_{\mathfrak c}=N^{\ell_{\mathfrak c}}\), \(z_{\mathfrak c}=N^{t_3}\) and
\(u_{\mathfrak c}=D_{\mathfrak c}^{\sigma^2}\). Let \(I_i\) be the actual
half-open exponent interval of coordinate \(i\) of the clipped cell. Then
\(\mathfrak c\subset I_1\times I_2\times I_3\). Put
\(\mathcal P_i=\{a\in\mathcal P_N:\log_Na\in I_i\}\). For every
\(a_i\in\mathcal P_i\), \(N^{t_i}\le a_i<N^{\rho t_i}\) and \(a_i<N^{\mu_*}\),
also for a clipped interval. Write \(t_i'\le\min(\rho t_i,\mu_*)\) for the
upper endpoint of \(I_i\). On \(\overline{\mathfrak c}\) the actual
coordinates satisfy \(x_1\ge x_2\ge x_3\ge u\). Two of the intervals are
equal or disjoint.

\paragraph{Exact occurrence map.}
By \eqref{eq:k2-aa-prefix-remainder} and Paper~I, (2.30), (2.32)--(2.33),
applied to \(\mathcal A_h\) at level \(D_{\mathfrak c}\), cutoff
\(z_{\mathfrak c}\le D_{\mathfrak c}^{1/2}\) and \(\epsilon_s=\sigma\),
\begin{equation}\label{eq:k2-r3-occurrence}
 r_3=\sum_{\mathfrak c\in\Xi^{\sharp}_3}\ \sum_{k\in\mathcal K_{\mathfrak c}}
       \sum_qO_{\mathfrak c,k}(q)\,\Delta_{A_b}(q,N),
 \qquad \mathcal K_{\mathfrak c}=\{\mathrm{R1}\}\cup\mathcal D^-(D_{\mathfrak c}),
\end{equation}
\[
 O_{\mathfrak c,k}(q)=\sum_{\substack{a_1>a_2>a_3\\a_i\in\mathcal P_i}}
 \sum_{\substack{d<D_{\mathfrak c}^\sigma\\d\mid P_N(u_{\mathfrak c})}}
 \lambda^{(k)}_d
 \sum_{\substack{(p_j)_{j\le\rho_k}\\ p_j\in\mathcal Q_{k,j}\\
 p_i\ne p_j\ ((i,j)\in I(\mathbf D))}}
 1_{q=a_1a_2a_3d\,p_1\cdots p_{\rho_k}} .
\]
Here \(\lambda^{(\mathrm{R1})}_d=\phi_d^-\) with no tail
(\(\rho_{\mathrm{R1}}=0\)). For \(k=\mathbf D\),
\(\lambda^{(k)}_d=\lambda_d^-(\mathbf D)\) and
\(\mathcal Q_{k,j}=\mathcal P_N\cap[D_j,\min(D_j^{1+\tau},z_{\mathfrak c}))\).
All coefficients have modulus at most one and depend only on
\((\mathfrak c,k)\). The identity is exact at every supported modulus.
Let \(M_{\mathfrak c}\in\{1,2,6\}\) be the number of orderings of the
equal intervals among \(I_1,I_2,I_3\). Let \(\Gamma_{\mathfrak c,k}\)
be \(M_{\mathfrak c}^{-1}\) times the same sum, with \(a_i\in\mathcal P_i\)
independent and no tail distinctness. Then
\begin{equation}\label{eq:k2-native13-coefficient}
 O_{\mathfrak c,k}=\Gamma_{\mathfrak c,k}-\Xi_{T,\mathfrak c,k}
 -\Xi_{P,\mathfrak c,k},
\end{equation}
where \(\Xi_T\) collects unrestricted terms with
\(p_i=p_j\) for some forbidden pair \((i,j)\in I(\mathbf D)\).
\(\Xi_P\) collects the remaining unrestricted terms with
\(a_i=a_j\) for an equal-interval pair \(I_i=I_j\). A prefix prime is
at least \(z_{\mathfrak c}\), and every tail prime is below it. The primes of
\(d\) are below \(u_{\mathfrak c}\le D_j\). Unequal grid intervals are
disjoint. Hence every modulus in the support of \(\Xi_T\) or \(\Xi_P\) is
divisible by \(p^2\) for a prime \(p\ge u_{\mathfrak c}\ge u_0\). Here
\(u_0=N^{\ell_0\sigma^2}\).

\paragraph{Support, cutoffs and constants.}
Every supported modulus is coprime to \(2N\). The cells are clipped
to \([u,\mu_*)\), and \(v_{\rm pk}(x)=(1+x)/2\) on \([0,\mu_*]\). Thus
\(\Theta_{\mathfrak c}\ge(1+t_1)/2\). The actual coordinate bound
\(x_1\ge x_3\ge u\) on \(\overline{\mathfrak c}\) also gives
\(\Theta_{\mathfrak c}\ge(1+u)/2\). Hence
\begin{equation}\label{eq:k2-r3-ell}
 \ell_{\mathfrak c}\ge\max\{A_3-t_3,2\rho t_3\}\ge\tfrac23A_3
 =\tfrac{33109}{240000},\qquad A_3=\tfrac12-\eta_L-\tfrac32\mu_*,
\end{equation}
and \(\ell_{\mathfrak c}-\ell_0\ge231463/24720000\). By the cubic tail
conditions, \(\sum_js_j<\ell_{\mathfrak c}\) for
\(s_j=\log_ND_j\). Thus every supported modulus satisfies
\begin{equation}\label{eq:k2-r3-support}
 q<N^{\rho(t_1+t_2+t_3)+\sigma\ell_{\mathfrak c}+\rho\ell_{\mathfrak c}}
 \le N^{e_Q},\qquad
 e_Q=(\rho+\sigma)(H_{\max}-\eta_L)<\tfrac35,\qquad H_{\max}=\tfrac{153}{256}.
\end{equation}
Since the cell ceiling gives \(0<\ell_{\mathfrak c}\le
H_{\max}=153/256<1\) and \(0<\sigma\le10^{-10}\),
\(\ell_{\mathfrak c}\sigma^2<\sigma^2<u=1/500\). Thus
\(u_{\mathfrak c}=N^{\ell_{\mathfrak c}\sigma^2}<N^u\). The cutoff
comparison uses only this exponent inequality.

Every labelled prime is at least \(u_{\mathfrak c}\), so the number of
labels is below \(3/(5\ell_0\sigma^2)\). With
\(M_\sigma=\lceil3/(5\ell_0\sigma^2)\rceil\) and the vector count
\(\exp(8\sigma^{-3})\), one has
\(|\Gamma_{\mathfrak c,k}|,|\Xi_{T}|,|\Xi_P|\le M_\sigma!\), and
\(\sum_{\mathfrak c,k}
 \bigl(|\Gamma_{\mathfrak c,k}|+|\Xi_T|+|\Xi_P|\bigr)
 \le C'_{\mathrm{K2},\sigma}\).

\paragraph{The squareful terms.}
Since \(|\Delta_{A_b}(q,N)|\le2(N/q+1)+2N/\varphi(q)\), Paper~I, (B.5),
with \(Y=u_0\) and \(Q_{\rm supp}=N^{e_Q}\) gives
\begin{equation}\label{eq:k2-r3-squareful}
 \sum_{\mathfrak c,k}\sum_q|\Xi_{T,\mathfrak c,k}(q)
 +\Xi_{P,\mathfrak c,k}(q)|\,|\Delta_{A_b}(q,N)|
 \ll_\sigma N^{1-\ell_0\sigma^2}L .
\end{equation}

\paragraph{Prime detection and the two Heath--Brown expansions.}
Put \(z_1=N^{1/(\log L)^3}\), so \(z_1\ge z_0(X)\) for every occurring
box scale \(\varepsilon N/2\le X\le N\).
Both source primes exceed \(z_1\). Apply Paper~I, (B.7), to \(1_{\mathbb P}(m)\)
and to \(1_{\mathbb P}(p')\). Apply Paper~I, (B.9), with roughness \(z_1\), to
\(\Lambda_{z_1}(m)\) with cap \(U_1=2N^{1/6}\) and to \(\Lambda_{z_1}(p')\)
with cap \(U_2=2N^{(1-b)/3}\). The identities apply since
\(m\le N^{1/2}<U_1^3\) and \(p'\le N^{1-b}<U_2^3\). For patterns
\(j,j'\in\{1,2,3\}\) the tested index is a labelled tuple
\(\mathbf v=(\mathbf v^{(1)},\mathbf v^{(2)})\in\mathcal R_{z_1}^{2j+2j'}\),
at most twelve coordinates, with
\[
 m=\textstyle\prod\mathbf v^{(1)},\quad p'=\prod\mathbf v^{(2)},\quad
 e=mp',\quad N^b\le m\le N^{1/2},\quad \varepsilon N<e\le N',
\]
the Möbius coordinates of \(\mathbf v^{(i)}\) at most \(U_i\), and coefficient
\(c_j(\mathbf v^{(1)})c_{j'}(\mathbf v^{(2)})/(\log m\log p')\).
Here \(c_j\) is the coefficient after (B.10). For a fixed product \(e\), put
\begin{equation}\label{eq:k2-pointmass-discrepancy}
 \mathcal S_e(q;N):=\Delta_{\delta_e}(q,N)
 =1_{e\equiv N\ (q)}-\frac{1_{(e,q)=1}}{\varphi(q)}.
\end{equation}
Thus the Heath--Brown tuple map uses the point-mass discrepancy
\(\mathcal S_e(q;N)\). The exact weighted point-mass expansion is
\[
 \Delta_{A_b}(q,N)=\sum_e A_b(e)\mathcal S_e(q;N).
\]
The weight \(A_b(e)\) keeps both eligible ordered orientations, as in
\eqref{eq:k2-source}. A prime square is counted once. The proper-power remainder
\(\mathcal E_{\rm pp}\) consists of the pairs in which \(m\) or \(p'\) is a
proper prime power. For every fixed \((\mathfrak c,k)\),
\begin{equation}\label{eq:k2-r3-hb-identity}
 \sum_q\Gamma_{\mathfrak c,k}(q)\Delta_{A_b}(q,N)
 =\mathcal E_{\rm pp}
 +\sum_{j,j'}\sum_{\mathbf v}
   \frac{c_j(\mathbf v^{(1)})c_{j'}(\mathbf v^{(2)})}{\log m\log p'}
   \sum_q\Gamma_{\mathfrak c,k}(q)\mathcal S_{e}(q;N),
\end{equation}
exactly. Each tuple keeps the natural discrepancy of its own product.
The per-cell coefficient bound is the factorial bound above; the sum over
cells uses the separate aggregate bound defined there.
Since \(e<N\), \(|c_jc_{j'}|/(\log m\log p')\le9\) and
\(\sum_q|\Gamma(q)||\mathcal S_e(q;N)|\le C_{\mathrm{K2},\sigma} F(e)\) with
\(F(e)=\tau(N-e)+2(1+L)\). The tuple multiplicity is at most
\(\tau_{12}(e)\). The divisor bounds \(\tau_{12}^2\le\tau_{144}\),
\(\tau^4\le\tau_{16}\) and Cauchy--Schwarz give
\begin{equation}\label{eq:k2-r3-moments}
 \sum_{e<N}\tau_{12}(e)F(e)\ll NL^{73},\qquad
 \sum_{e<N}\tau_{12}(e)F(e)^2\ll NL^{79}.
\end{equation}
There are \(\ll N^{1-b/2}L\) proper-power pairs, and
\(\sum_{e<N}\tau(e)F(e)^2\ll NL^9\). Therefore
\begin{equation}\label{eq:k2-r3-pp}
 |\mathcal E_{\rm pp}|\ll_\sigma(N^{1-b/2}L)^{1/2}(NL^9)^{1/2}
 =N^{1-b/4}L^5 .
\end{equation}

\paragraph{Boxes, logarithm freezing and faces (H6).}
Put \(\delta=L^{-334}\). Use the equal-ratio mesh of Paper~I, (B.12),
with \(\delta_X\) replaced by \(\delta\), on every ambient coordinate axis.
Keep \(\{1\}\) as an atom. A box \(\mathcal B\) is a Cartesian product
of cells. Its lower endpoints are \(V_i\), and \(X=\prod_iV_i\). Let
\(\lambda_1,\lambda_2\) be the sums of \(\log V_i\) over the two sides.
Then \(\varepsilon N/2\le X\le N\) on every box meeting the source, so
\(K_{\rm scale}=1\). Replacing \(1/(\log m\log p')\) by
\(1/(\lambda_1\lambda_2)\) changes the weight by \(O(L^{-335})\), since
\(\lambda_1\ge bL/2\). By \eqref{eq:k2-r3-moments},
\begin{equation}\label{eq:k2-r3-log}
 |E_{\rm LOG}|\ll_\sigma NL^{72-334}.
\end{equation}
The faces are \(m=N^b\), \(m=N^{1/2}\), \(e=\varepsilon N\), \(e=N'\), and
the cap faces \(v_i=U_1\), \(v_i=U_2\) on Möbius coordinates. All are
singletons. On each, an active coordinate is at least \(N^{1/20}\).
Paper~I, Lemma~(B.14), with \(s\le12\), \(K_\gamma=1\) and
\(\eta_{\rm dist}\le12\delta\), gives at most \(\ll NL^{11-334}\) boundary
tuples. Cauchy--Schwarz with \eqref{eq:k2-r3-moments} gives
\begin{equation}\label{eq:k2-r3-bdry}
 |E_{\rm bdry}|\ll_\sigma NL^{45-334/2}.
\end{equation}
Thus (H6) holds with \(c_{\log}=72\), \(c_R=45\). On a nonboundary box every
source cut and cap is constant. The coefficient is the outer scalar
\(c_{\mathcal B}=(-1)^{j+j'}\binom3j\binom3{j'}(\log X)^2/(\lambda_1\lambda_2)\),
times
\(w(\mathbf v)=\prod_iw_i(v_i)\) with
\(w_i\in\{\mu_{z_1},1_{\mathcal R_{z_1}},(\log v/\log X)1_{\mathcal R_{z_1}}\}\).
Each side carries exactly one logarithmic coordinate, and
\((\log X)^2/(\lambda_1\lambda_2)
\le 1/(b(1-b))+o(1)\le4.5+o(1)\), hence the scalar norm bound
\(\lVert c_{\mathcal B}\rVert_\infty\le45\).
Each coordinate meets at most \(2L^{335}\) cells, so
\begin{equation}\label{eq:k2-total-coefficient-norm}
 \sum_{\mathfrak c,k}\sum_{j,j'}\sum_{\mathcal B}|c_{\mathcal B}|
 \ll_\sigma L^{4020}\le L^{4022},\qquad K_0=4022\le24C_{\rm box}+26=8042.
\end{equation}

\paragraph{Classification.}
Let \(e_i=\log_NV_i\), \(\mathcal W=[60\sigma,1/3+14\sigma]\), and
\(\kappa=(\log2+12\log(1+\delta))/L\le\sigma/202\). A box is Type~II if
some \(e_i\in\mathcal W\). Select the largest such \(e_i\), breaking ties
by the fixed coordinate order. Otherwise every Möbius coordinate is small, since
\(1/6+\kappa\) and \((1-b)/3+\kappa\) are below \(1/3+14\sigma\). On the
\(m\)-side, \(\sum e_i\ge b-\kappa>360\sigma\). On the \(p'\)-side,
\(\sum e_i\ge1/2-\kappa-\log(1/\varepsilon)/L>360\sigma\). These are the only side bounds
needed to rule out an all-small side. Together with the global Heath--Brown
sum bound \(\sum e_i<1+\kappa\), the threshold \(1/3+14\sigma\) permits
at most two large variables. Hence each side has exactly one large smooth
coordinate. Every other box has exactly two large smooth coordinates \(N_1\ge N_2\ge X^{1/3}\). It has at most ten small
coordinates, of product \(M\), with
\begin{equation}\label{eq:k2-r3-typeI2}
 \log_XM\le\frac{(600+1/202)\sigma}{1-\sigma/101}<6.02\cdot10^{-8},
 \qquad N_1N_2M\le2^{30}X,\qquad q<X^{0.5977}.
\end{equation}
There are no Type-I-1 boxes.

\paragraph{Type-II boxes: Paper~I, (H1)--(H5) and LEN.}
Fix \((\mathfrak c,k)\) and a Type-II box. Let its selected lower endpoint
be \(T\), put \(\xi=\log_NT\), and let the other lower endpoints have product \(M=X/T\).
\begin{description}
\item[H1.] \(C_{\mathcal B}=\Gamma_{\mathfrak c,k}\) has the form (3.4). Here
\(J\) is the three prefix labels and the tail labels. The prime sets are
\(\mathcal P_i\) and \(\mathcal Q_{k,j}\), \(u_L=u_{\mathfrak c}\ge u_0\),
and \(d\mid P_N(u_{\mathfrak c})\) is squarefree. Repeated primes are kept.
We have \(q\le X^{0.6}\) by \eqref{eq:k2-r3-support} once
\((3/5-e_Q)L\ge(3/5)\log(2/\varepsilon)\). Also \((q,2N)=1\). The
selector \(\chi_{\mathfrak c}\) is constant on the block. The removed
corrections are \(\Xi_T,\Xi_P\) and \(\mathcal E_{\rm pp}\). There is no
comparator correction, since the comparator is natural on every tuple.
\item[H2.] Take the upper exponents \(t_1'\ge t_2'\ge t_3'\) and
\(\min(\rho s_j,t_3)\). For \(d_*=\rho(\Theta_{\mathfrak c}-\eta_L)\), the
merged list satisfies \(\mathrm{WF}(d_*)\). The reserves are
\(\rho(1/2-\eta_L-3\mu_*/2)=33109\rho/160000\),
\(\rho(1/2-\eta_L-5\mu_*/2)=1859\rho/160000\), \(\rho(\ell_{\mathfrak c}-t_3)
\ge\rho(2\rho-1)t_3\). For a tail rank, the reserve is \(\rho\) times the strict
inequality \(\sum_{i<j}s_i+2s_j<\ell_{\mathfrak c}\) from the
\(\mathcal D^-\) conditions. Put \(h=d_*+60\sigma\), \(\delta_{\rm pk}=30\sigma\),
\(g=\xi-30\sigma\). Then \(h-1/2\ge9\cdot10^{-4}\),
\(\Theta_{\mathfrak c}-h>0.0000999939\), \(h\le H_{\max}\),
\[
30\sigma\le g\le1/3-16\sigma,
\qquad \delta_{\rm pk}=30\sigma,
\qquad g\le1/3-\delta_{\rm pk}/2.
\]
The upper requirement has reserve at least \(\sigma\), since
\((1/3-15\sigma)-(1/3-16\sigma)=\sigma\). Also
\(h\le\Theta_{\mathfrak c}\le\widetilde H(t_1',t_2',t_3')\) because
\((t_1',t_2',t_3')\in\overline{\mathfrak c}\) and \(t_1'\le\mu_*\).
Paper~I, Theorem~4.1, gives \(J=J_1\sqcup J_2\sqcup J_3\). Put \(d\) in
slot one, with \(\gamma_s\) as in Paper~I, (3.5). The tuple map gives
\(\Gamma_{\mathfrak c,k}=M_{\mathfrak c}^{-1}C_{\rm slot}(\gamma_1*\gamma_2*\gamma_3)\),
\(C_{\rm slot}\le M_\sigma!\), \(|\gamma_s|\le1\). Slots two and three are
\(u_0\)-rough, and empty slots carry the atom \(1_{v=1}\). The scalar
\(M_{\mathfrak c}^{-1}\) is an outer coefficient.
\item[H3.] Let \(Q,R,S\) be the scales of \(\gamma_1,\gamma_3,\gamma_2\). The
loads (4.3), the slot-one load \(\sigma\ell_{\mathfrak c}\le\sigma H_{\max}\)
of \(d\), and \(A_0=X^{1+\sigma/20}\) give reductions
\((7527/256,\,60,\,114421/1280,\,26247/256)\sigma\). They are, in turn, in
\(Q/T\), \(T^2R^2S\), \(A_0^{7/32}TQR^2S^5\), \((TR/S)^{7/16}TQR^2S^5\).
The required reductions are \(29,8,15,15\), with losses at most \(2\omega\),
\(\omega=\log_N(N/X)\le\sigma/202\). We also have \(\log_N(QRS)\le e_Q\), so
\(QRS<M_IX^{-89\sigma}\) with reserve \(>0.0691\). Next, \(N_I\ge QX^{\sigma/3}\)
with reserve \(29.069\sigma\), and \(M_I\ge X^{1/2}\). Thus (3.6)--(3.7)
hold. The selected coefficient is
\(\alpha_n=w_{\rm sel}(n)1_{I_T}(n)\in\{\mu(n),1,\log n/\log X\}1_{P^-(n)\ge z_1}1_{I_T}(n)\),
with \(\zeta=z_1\ge z_0(X)\) and \(I_T=(T,(1+\Delta)T]\). The complementary
coefficient \(\beta_{m^\sharp}\), the product of the other at most eleven
weights, satisfies \(|\beta_{m^\sharp}|\le\tau_{11}(m^\sharp)\le\tau_{23}(m^\sharp)\).
Here \(\tau_{23}=d_{23}\) is the 23-fold divisor function, with
\(d_{23}(n)=\sum_{n=n_1\cdots n_{23}}1\).
As in Paper~I, (B.22)--(B.23), with at most eleven complementary cells,
\(N_I\ge T/(6L^{334})\), \(M_I=M((1+\Delta)^{r}-1)\asymp ML^{-334}\),
\(m_0=M(1+(1+\Delta)^r)/2\le2X/T\), so \(C_{\rm ctr}=2\). On the doubled
support, \(mn\le X(1+\Delta)(1+\tfrac32((1+\Delta)^{11}-1))\le N'(1+20\delta)\),
because \(X\le e\le N'\) for every tested product of the box. Hence
\(N-mn\ge NL^{-334}(L-20)\ge1\) and \(|mn-N|\le N<X^3\), which is (3.10).
The integer convention of \(I_T\) is that of Paper~I, Lemma~B.22, since the
mesh and the roughness are the same. LEN follows from the condition
\begin{equation}\label{eq:k2-med-length}
 \sigma L\ge2^{130}(335)(\log L+1).
\end{equation}
It gives \(T^{2^{-126}}\ge12L^{334}\) and
\(N_I\ge(2T)^{1-2^{-126}}\).
\item[H4.] \(a=N\), \(0<N\le XL\le A_0\), and every supported modulus is
coprime to \(2N\).
\item[H5.] \(N^{60\sigma}\le T\le N^{1/3+14\sigma}\), and (3.99) gives
\(T\le X^{1/3+15\sigma}\) and \(T<X^{1/2}\).
\end{description}
The selected-factor Siegel--Walfisz bound is Paper~I, (3.25). It is
proved for these three coefficient classes with any \(\zeta\ge z_0\). In the
large-norm case of Paper~I, (3.48), the zero mode uses
Lemma~\ref{hyp:ra-ms-transfer}(ii) in place of \cite[Theorem~0(a)]{BFI86}.
Its premises (3.1), (3.2), (3.6), (3.8) with LEN, (3.11) and (3.25)--(3.27)
are those verified above. Theorem~2.8 of Paper~I enters only inside the
proof of Theorem~3.2, with the dictionary
\begin{equation}\label{eq:k2-theorem-28-dictionary}
 r_T=z,\quad s_T=p_0,\quad n_T=p_0m,\quad a_T=\varepsilon_yN,
 \quad N_T=8p_0H_*SF,\quad A_T=A_0,
\end{equation}
both signs, all nine smooth weights and the full bracket.
Paper~I, Theorem~3.2, (3.13), gives, uniformly in the box, the block and
the cell,
\begin{equation}\label{eq:k2-med-per-box}
 |S_{\mathcal B}^{\rm PSDT}|
 \ll_{\sigma,A'}XL^{-A'}+X^{1-c_\sigma}L^C,\qquad
 S_{\mathcal B}^{\rm PSDT}=\sum_q\Gamma_{\mathfrak c,k}(q)
 \sum_{\mathbf v\in\mathcal B}w(\mathbf v)\mathcal S_{e}(q;N).
\end{equation}
With \(A=4\), \(C_{\rm box}=334\), \(A'=A+24C_{\rm box}+31=8051\) and
\eqref{eq:k2-total-coefficient-norm}, Paper~I, (3.92), gives
\begin{equation}\label{eq:k2-med-aggregate}
 E_{\rm II}=\sum|c_{\mathcal B}|\,|S_{\mathcal B}^{\rm PSDT}|
 \ll_{b,\sigma}NL^{4022-8051}+N^{1-c_\sigma}L^{4022+C}.
\end{equation}
Here \(S_{\mathcal B}^{\rm PSDT}\) already contains
\(\Gamma_{\mathfrak c,k}\), hence its factor
\(M_{\mathfrak c}^{-1}\) from \eqref{eq:k2-native13-coefficient}. The
coefficient norm \eqref{eq:k2-total-coefficient-norm} sums only
\(|c_{\mathcal B}|\), so the orbit factor is paid exactly once.

\begin{lemma}[Relaxed Type-I-2 estimate]\label{lem:k2-I2-prime}
The source hypotheses are the Type-I-2
conditions (SH1)--(SH5) stated in Paper~I:
\begin{quote}
For a direct shell \([X,2X)\), assume the following source conditions
(SH1)--(SH5).
\begin{itemize}
\item[(SH1)] The coefficient \(\gamma(q)\) is independent of tested variables,
\((q,2N)=1\), \(|\gamma(q)|\le C_\sigma\), \(q\le N^{e_Q}\),
\(e_Q<153/256\), and \((a,q)=1\) on its support.
\item[(SH2)] The tested geometric source and cap restrictions are constant
on each nonboundary box; the roughness sieve is separate as in (SH3).
Smooth factor weights have \(\sup|w|+V(w)\ll1\).
\item[(SH3)] All HB factors are rough at the threshold \(z^*=z_0(X)\) or
\(z_0(N)\).
\item[(SH4)] The discrepancy is exactly (3.3).
\item[(SH5)] The assignment priority is proper prime powers, boundary boxes,
Type II, Type I-1, then Type I-2.
\end{itemize}
\end{quote}
Assume (SH1)--(SH5) for a nonboundary K2 box. Let the two large smooth
variables have nonnegative weights \(w_1,w_2\) and \(N_1,N_2\ge X^{1/3}\).
Let the remaining small variables, at most ten, have product \(M\). Assume
\[
 M\le X^{6.02\cdot10^{-8}},\qquad N_1N_2M\le2^{30}X,
 \qquad q<X^{0.5977},
\]
and \(|\alpha_m|\ll\tau_{10}(m)\le\tau(m)^9\). Then, for every fixed
\(A>0\), uniformly in the supported unit residues,
\[
\sum_{q\asymp Q}|\gamma(q)|\left|\sum_m\alpha_m
 \sum_{n_i\in I_i,\ P^-(n_i)\ge z^*}
 w_1(n_1)w_2(n_2)\mathcal S_{mn_1n_2}(q;a)\right|
\ll_{\sigma,A} XL^{-A}.
\]
The same statement holds for a fixed finite sum of signed weights. Split
each bounded real weight into its two nonnegative parts.
\end{lemma}

\begin{proof}
Restrict first to \((m,q)=1\), and write \(\Delta=C-M_q\) for the positive
congruence and comparator functionals. In the K2 application the inner
product is \(e=mn_1n_2\), the residue is \(N\), and \((q,N)=1\). If
\((m,q)>1\), then \((e,q)>1\), so both \(1_{e\equiv N\ (q)}\) and the
comparator \(1_{(e,q)=1}/\varphi(q)\) vanish. The omitted range contributes
zero. Apply the Brun--Bonferroni sandwich
of Paper~I, (3.102)--(3.104), at
\(D_f=X^{\epsilon_I/2}\), \(\epsilon_I=7.5\sigma\). For a requested
fixed logarithmic precision \(A_{\rm sv}\), put
\[
 s_f=\frac{\log D_f}{\log z^*},\quad
 E=\sum_{p<z^*}\frac1p,\quad J=2\lfloor s_f/2\rfloor,
\]
and increase the common threshold until
\(J\ge\max\{2,e^2E,A_{\rm sv}\log L\}\). Since
\(s_f\asymp\sigma(\log\log N)^3\) and \(E=O(\log\log N)\), this is
eventual at fixed \(\sigma\), and \(E^J/J!\le L^{-A_{\rm sv}}\).

For the two large factors put \(R_i=1_{P^-(n_i)\ge z^*}\) and
\(G_i=\Psi_i^+-\Psi_i^-\ge0\). Positivity gives
\[
 \Psi_1^+\Psi_2^+-G_1\Psi_2^+-\Psi_1^+G_2
 \le R_1R_2\le\Psi_1^+\Psi_2^+,
\]
so the three signed gaps and the two nonnegative comparator gaps
bound \(|\Delta(R_1R_2)|\), as in Paper~I, (3.113). Expand the
three gaps over \(d_i\le D_f\). The completed hyperbola estimate (3.109)
 gives \(O_{\epsilon_X}(D_f^2q^{1/2}X^{\epsilon_X})\). Here we fix
\(\epsilon_X=10^{-4}\). Squareful moduli are included
by the full-modulus Weil--Estermann bound.

The two one-dimensional sieve gaps satisfy
\[
 \sum_{n\in I_i}G_i(n)w_i(n)\ll_{A_{\rm sv}}N_iL^{-A_{\rm sv}}+D_f,
 \qquad \sum_{n\in I_i}\Psi_i^+(n)w_i(n)\ll N_i+D_f.
\]
Thus the total is bounded by
\[
 M D_f^2Q^{3/2}X^{\epsilon_X}L^C+XL^{C-A_{\rm sv}}
 +M D_f(N_1+N_2)L^C+M D_f^2L^C. \tag{I-2'}
\]
The ten-variable check uses the following upper bound. The expression is increasing
throughout \(0<\sigma\le10^{-10}\), so its maximum is evaluated at
\(\sigma=10^{-10}\):
\[
 \frac{600\sigma+\kappa}{1-\omega}
 \le\frac{(600+1/202)\sigma}{1-\sigma/101}
 \le\frac{(600+1/202)10^{-10}}{1-10^{-10}/101}
 =\frac{121201}{2019999999998}
 \approx6.0000495049564356926\cdot10^{-8}<6.02\cdot10^{-8},
\]
where \(\omega,\kappa\le\sigma/202\) absorb all fixed dyadic factors,
so \(1-\omega\ge1-\sigma/202\ge1-\sigma/101\), consistently with
\eqref{eq:k2-r3-typeI2}.
Eleven small variables would give
\(\frac{133321}{2019999999998}\approx
6.600049504957030\cdot10^{-8}\), which exceeds the allowed bound.
With \(q<X^{0.5977}\), the first exponent in (I-2') is at most
\[
 6.02\cdot10^{-8}+7.5\cdot10^{-10}+1.5(0.5977)+\epsilon_X
 =0.89665006095<0.897,
\]
leaving slack \(0.00034993905\). Also, using
\(N_1N_2M\le2^{30}X\), \(N_1\ge N_2\ge X^{1/3}\), and
\(D_f=X^{3.75\sigma}\),
\[
 M D_f(N_1+N_2)\le2^{31}D_fX/N_2
 \ll X^{2/3+3.75\cdot10^{-10}}<X^{0.668},
 \qquad
 M D_f^2\le X^{6.02\cdot10^{-8}+7.5\cdot10^{-10}}<X^{0.668}.
\]
The coefficient bound uses
\(\tau_{10}(m)\le\tau(m)^9\), so \(B_0=9\) is fixed. Finally, a family of
at most \(O_\sigma(L^K)\) boxes is absorbed by requesting
\(A_{\rm sv}\ge K+A+5\), after all fixed logarithmic losses. This proves
the lemma.
\end{proof}

\paragraph{Type-I-2 boxes: the complementary estimate.}
Lemma~\ref{lem:k2-I2-prime} is a relaxed form of Paper~I,
(3.111), for this source. Below we check each displayed hypothesis
against the actual K2 source.

For the present K2 Type-I-2 box, the nontrivial SH2 check is the
display
\begin{equation}\label{eq:k2-sh2-check}
\begin{gathered}
 1\le v\le X,
 \qquad w(v)\in\left\{1,\frac{\log v}{\log X}\right\},
 \qquad P^-(n_i)\ge z^*\ \text{is a separate tested condition},\\
 \sup|1|+V(1)=1,\qquad
 \sup\left|\frac{\log v}{\log X}\right|
   +V\!\left(\frac{\log v}{\log X}\right)
 \le 1+\frac{2\log(1+L^{-334})}{\log X}=O(1).
\end{gathered}
\end{equation}
Thus the two tested large factors have the required nonnegative weights and
bounded supremum plus variation. The roughness condition is not hidden inside
the weight.

For the Type-I-2 boxes above, (SH1) has
\(\gamma=\Gamma_{\mathfrak c,k}\). The support and bound are
\eqref{eq:k2-r3-support} and \eqref{eq:k2-native13-coefficient},
\(e_Q<153/256\), and the residue is the unit \(a=N\). The partition
construction proves the geometric constancy in (SH2) on each nonboundary
box. In this two-large-variable case the two tested large factors are the
selected smooth factors. The roughness restriction is kept as the separate
condition \(P^-(n_i)\ge z^*\) in the inner sum. Their smooth weights are then
individually \(1\) or \(\log v/\log X\), not \(\mu_{z_1}\). Hence they are
nonnegative with \(\sup|w|+V(w)\ll1\) as required by (SH2). The
\(\mu_{z_1}\) factor occurs only among the complementary small factors.
Condition (SH3) is \(z_1\ge z_0(X)\).
Condition (SH4) is the natural full-product discrepancy fixed above.
Condition (SH5) is the priority in the preceding boundary and Type-II
classification.

On a large-coordinate cell \((V,(1+\delta)V]\),
\(\delta=L^{-334}\), one has \(1\le V\le X\). The smooth weight
\(w(v)=1\) has supremum plus variation equal to 1. The alternative
\(w(v)=\log v/\log X\) is nonnegative and increasing, with
\[
 \sup |w|+V(w)
 \le 1+\frac{2\log(1+\delta)}{\log X}=O(1).
\]
The roughness restriction remains a separate tested condition under (SH3)
and inside the sieve, and it is not claimed to be constant on a box.

Thus the relaxed estimate of this chapter applies with
\(z^*=z_1\), the two large smooth weights nonnegative, and
\[
 |\alpha_{m^\sharp}|\le\tau_{10}(m^\sharp)\le\tau(m^\sharp)^9.
\]
Each of the at most ten small labels contributes one factor. Normalized
logarithmic weights \(\log v/\log X\) are at most one on the support.
Thus the coefficient hypothesis uses the fixed divisor exponent
\(B_0=9\), independent of \(N\). The first term of (I-2') uses the
following values. The slack is \(0.00034993905\)
at exponent \(0.89665006095\). The two final terms have the bound
\(2/3+3.75\cdot10^{-10}\). The absorption is
\(\omega,\kappa\le\sigma/202\), and the divisor bound is \(\tau_{10}\)
with \(B_0=9\). By \eqref{eq:k2-r3-typeI2}, Lemma~\ref{lem:k2-I2-prime}
applies to every K2 Type-I-2 box. Requesting \(A_I=K+A+5\) for the
fixed \(O_\sigma(L^K)\) box family (here \(K=4022\)) gives
\begin{equation}\label{eq:k2-r3-complementary}
 E_{\rm I2}\ll_{b,\sigma}NL^{-A}=NL^{-4}.
\end{equation}

\begin{proposition}[Row-three remainder]\label{prop:k2-row3}
Fix \(b\in(1/3,1/2)\) and \(0<\sigma\le10^{-10}\). Under Paper~I,
Theorems~2.8 and~3.2, with the relaxed Type-I-2
Lemma~\ref{lem:k2-I2-prime} and Lemma~\ref{hyp:ra-ms-transfer}(ii) at (3.48),
\[
 \frac{|r_3|}{\mathcal N_K}\ll_{b,\sigma}
 N^{-\ell_0\sigma^2}L^3+N^{-b/4}L^7+L^{-260}+L^{-120}+L^{-2}
 +N^{-c_\sigma}L^{4024+C}.
\]
\end{proposition}

\begin{proof}
By \eqref{eq:k2-r3-occurrence} and \eqref{eq:k2-native13-coefficient},
\(r_3\) is the \(\Gamma\)-sum minus the squareful sum
\eqref{eq:k2-r3-squareful}. By \eqref{eq:k2-r3-hb-identity}, the
\(\Gamma\)-sum is \(\mathcal E_{\rm pp}\) plus the tuple sum. Every tuple
lies in exactly one box. The boundary tuples give
\eqref{eq:k2-r3-bdry}. On the other boxes logarithm freezing costs
\eqref{eq:k2-r3-log}. Each box is of Type~II or Type~I-2. Sum
\eqref{eq:k2-r3-squareful}, \eqref{eq:k2-r3-pp}, \eqref{eq:k2-r3-log},
\eqref{eq:k2-r3-bdry}, \eqref{eq:k2-med-aggregate} and
\eqref{eq:k2-r3-complementary}, with \(O_\sigma(1)\) pairs
\((\mathfrak c,k)\). Dividing by
\(\mathcal N_K=(1-\varepsilon)C(N)N/L^2\ge(1-\varepsilon)C_2N/L^2\gg_b N/L^2\) gives the
claim.
\end{proof}

\begin{proposition}[Exterior source estimates]\label{prop:k2-applications}
Fix \(b,\sigma\) before \(N\). For the source \eqref{eq:k2-source}:
\begin{enumerate}
 \item \eqref{eq:k2-parent-comparison} and \eqref{eq:k2-main-mass} hold
 for the two-orientation source and four signed AA main terms in
 \eqref{eq:k2-row-main-masses} and \eqref{eq:k2-parent-functional};
 \item \eqref{eq:k2-aa-density-mass} and \eqref{eq:k2-aa-mass} hold for
 \(N\ge N_{\mathrm{AA}}(b,\sigma)\);
 \item \eqref{eq:k2-peel} holds for \(N\ge N_{\rm peel}(\sigma)\);
 \item \eqref{eq:k2-remainder-estimates} holds.
\end{enumerate}
\end{proposition}

\begin{proof}
Item 1 is \eqref{eq:k2-parent-reserve}--\eqref{eq:k2-parent-kernel} and the
derivation after \eqref{eq:k2-main-mass}. Item 2 is
\eqref{eq:k2-aa-decay}--\eqref{eq:k2-aa-density-derived}. Item 3 is
\eqref{eq:k2-peel-ranges}--\eqref{eq:k2-peel}. Item 4 is
\eqref{eq:k2-low-row-remainders} and Proposition~\ref{prop:k2-row3}.
\end{proof}


\begin{proposition}[Exterior comparison]\label{prop:k2-comparison}
For fixed \(b,\sigma\) there is a nonnegative \(r_K(N;b,\sigma)\to0\) such
that, for \(N\ge T_r(b,\sigma)\) with \(T_r\) as defined below,
\begin{equation}\label{eq:k2-final-fixed-sigma}
\begin{aligned}
 \frac{4B_b(N)}{W(N)}
 &\le c_{K2,G0}I(b)+A_K(b)C_{\mathrm{AA}}\sigma
          +P_K\sigma^9+r_K(N;b,\sigma),\\
 P_K&=\frac{388400000}{1-\varepsilon}
     =\frac{3884000000000}{9999}.
\end{aligned}
\end{equation}
Here
\(T_K=\max\{N_{\mathrm{mass}}(b,\sigma),N_{\mathrm{AA}}(b,\sigma),
N_{\rm peel}(\sigma),T_{\mathrm{len}}\}\). The threshold \(N_{\mathrm{mass}}\) is from
the fixed-parameter estimate at \eqref{eq:k2-main-mass}.
\(N_{\mathrm{AA}}\) is from
item~2 of Proposition~\ref{prop:k2-applications}. Both exist with
\(b,\sigma\) fixed. Finally, \(T_{\mathrm{len}}\) is the named length and
application threshold of Section~\ref{sec:notation}. In addition,
\begin{equation}\label{eq:k2-error-definition}
 r_K=\frac{E_{\rm top}+[M_K-\mathcal N_K2e^{-\gamma}I(b)G_*]_+
 +\sum_{j=0}^3|r_j|}{\mathcal N_K}
 =o_{b,\sigma}(1).
\end{equation}
In particular, let \(\widehat T_r\) be the least even threshold above
which \(r_K\le\delta_{\rm err}/10\), and set
\(T_r=\max\{N_{\mathrm{mass}}(b,\sigma),N_{\mathrm{AA}}(b,\sigma),
T_K,\widehat T_r\}\).
These thresholds exist at fixed \(b,\sigma\).
\end{proposition}

\begin{remark}[Single payment of each error]
In \cite[Theorem~2.6, (2.30)]{PaperI} the three summands
\(XV E_s\), \(R_1^{\nu}\) and \(R^{\nu}\) are distinct. The first is
paid only by \eqref{eq:k2-aa-mass}, the other two only by
\eqref{eq:k2-remainder-estimates}. The decaying part of
\eqref{eq:k2-aa-decay} and the \(o(1)\) in
\eqref{eq:k2-aa-density-derived} are absorbed by thresholds at fixed
\(\sigma\). The peel term \(P_3\) contains no linear-sieve
application. Thus
\[
 r_K=\frac{E_{\mathrm{top}}+\sum_j|r_j|
       +[M_K-\mathcal N_K2e^{-\gamma}I(b)G_*]_+}{\mathcal N_K},
\]
and no error is counted twice.
\end{remark}

\begin{proof}
Use \eqref{eq:k2-signed-join} and divide by \(\mathcal N_K=W/4\). The
principal term is bounded by \(\mathcal N_K2e^{-\gamma}I(b)G_*\) plus its
positive part. Also \(2e^{-\gamma}G_*\le c_{K2,G0}\) by
\eqref{eq:k2-certificate}--\eqref{eq:k2-gamma-enclosure}.
Proposition~\ref{prop:k2-applications}(2),(3) pay the sieve and peel terms
above \(T_r\). By definition, \(T_r\) includes
\(N_{\mathrm{mass}}(b,\sigma)\), \(N_{\mathrm{AA}}(b,\sigma)\) and
\(T_K\). The remaining terms form \(r_K\). First,
\(E_{\rm top}/\mathcal N_K\ll L^{-331}\). The positive part contributes
\(o_{b,\sigma}(1)\) by \eqref{eq:k2-main-mass}. The row remainders are
\(O_{b,\sigma}(L^{-2})\) by \eqref{eq:k2-remainder-estimates}. Hence
\(r_K=o_{b,\sigma}(1)\).
The convergence is ineffective through Siegel--Walfisz, so \(T_r\)
exists but is not computed.
\end{proof}
